\documentclass{rvcepaper}
\begin{document}

% ================= BLOCK A: 2025-26 odd-semester paper (inserted after the Apr-Jun 2026 block)
% ---------- HS251TA Principles of Management and Economics, Improvement CIE, 12.01.2026
\begin{iempaper}
\paperinfo{shownote=0,date={12.01.2026},exam={Improvement CIE},maxmarks={10 + 50},sem={V},prog={UG},duration={30 + 90 Min},
 title={Principles of Management and Economics},code={HS251TA}}
\iemhead
{\small Note:}\par\vspace{-1pt}
\hspace*{0.5cm}{\footnotesize 1. \hspace{0.1cm}Answer all the Questions.}\par\vspace{3pt}
\begin{cietable}{M}{BT}{CO}
\qpart{Part -- A}
\q{1}{What is a need in the motivation process?}{2}{2}{3}
\q{2}{State Maslow's hierarchy of needs.}{2}{1}{3}
\q{3}{Mention any two characteristics of leadership.}{2}{1}{3}
\q{4}{What is time management?}{2}{2}{2}
\q{5}{Mention any two stress management techniques.}{2}{1}{2}
\qpart{Part -- B}
\q{1}{During an economic slowdown, the central bank reduced interest rates and injected liquidity into the banking system to revive economic activity.\newline
a) Explain the role of money and banks in economic recovery.\newline
b) Describe how commercial banks influence the money supply.}{10}{3}{4}
\q{2}{What are the different classifications of unemployment? With a neat diagram explain how it is related to GDP?}{10}{3}{4}
\q{3}{Explain Keynesian model of Macroeconomics and how it is different from classic model in macroeconomics.}{10}{2}{4}
\q{4}{Explain Abraham's Maslow's theory of human motivation with hierarchy diagram.}{10}{2}{3}
\q{5}{What is transactional and transformational leadership? Explain the futures of each type of leadership.}{10}{3}{3}
\end{cietable}
\end{iempaper}

% ================= BLOCK B: Re-CIE of 30/06/2025 (inserted right after the 25/06/2025 Mathematical Modelling page)
% ---------- MA266TEU Mathematical Modelling, Re-CIE, 30/06/2025
\begin{xpaper}
\paperinfo{deptline={Department of Mathematics},date={30/06/2025},exam={Re-CIE},maxmarks={100},sem={VI},prog={UG},duration={2 Hrs. (2pm-4pm)},
 title={Mathematical Modelling},code={MA266TEU}}
\xhead
\begin{xtable}{Quiz: Questions}
\q{1}{The angle of projectile to attain the maximum horizontal distance should be \blank[1.2cm] the maximum height should be \blank[1.2cm].}{2}{L1}{2}
\q{2}{A hot object at $120^\circ$C is placed in a room with temperature $30^\circ$C. If the cooling constant is $k = 0$, what is the temperature after 5 minutes?}{2}{L1}{1}
\q{3}{In the long run, the proportion of time that the given Markov chain will be in state A is \blank[1.2cm] and state B is \blank[1.2cm].
\par\vspace{2pt}\centerline{\begin{tikzpicture}[x=1cm,y=1cm,font=\small,>=Stealth,state/.style={draw,circle,minimum size=0.7cm,inner sep=0pt}]
\node[state] (A) at (0,0) {A}; \node[state] (B) at (3,0) {B};
\draw[->] (B) to[bend right=35] node[above]{0.8} (A);
\draw[->] (A) to[bend right=35] node[below]{0.7} (B);
\draw[->] (A) to[out=150,in=210,looseness=8] node[left]{0.3} (A);
\draw[->] (B) to[out=30,in=-30,looseness=8] node[right]{0.2} (B);
\end{tikzpicture}}\par\vspace{16pt}}{2}{L2}{3}
\q{4}{Define probability vector. Check whether the vector $\left(\frac{1}{2},\frac{1}{4},\frac{1}{4},0\right)$ is a probability vector or not.}{2}{L2}{2}
\q{5}{The growth of an organism follows the law $\dfrac{dx}{dt} = \dfrac{1}{13}x$. If initially there were 100 of them, then after 13 years, their approximate number will be \blank[1.2cm] and their doubling period is approximately \blank[1.2cm] years.}{2}{L1}{1}
\end{xtable}
\ptitle{Test}
\begin{longtable}{|C{1.05cm}|p{12.2cm}|C{0.95cm}|C{0.95cm}|C{0.95cm}|}
\hline \textbf{SL. No} & \multicolumn{1}{c|}{\textbf{Test: Questions}} & \textbf{M} & \textbf{BT} & \textbf{CO}\\\hline
\q{1.}{A drug is induced in a patient's bloodstream at regular intervals of duration $T$ each. Simultaneously the drug disappears at a rate proportional to the concentration $c(t)$ of the drug present at any time $t$. Determine the differential equation governing $c(t)$, solve it and interpret the results.}{10}{L3}{2}
\q{2.}{Write the systematic 12-step procedure for tackling a problem through the use of mathematical modelling.}{10}{L1}{1}
\q{3a.}{Describe the concept of a simple compartment model with an example.}{05}{L2}{2}
\q{3b.}{The territory of a sales man consists 3 cities. A, B, C. he never sells in the same city on successive days. If he sells in city A, then the next day he sells in B. However, if he sells either in B or C, then the next day he is twice likely to sell in city A as in the other city. How often does he sell in each of the cities in steady state?}{05}{L2}{3}
\q{4a.}{Using the limiting behaviour of the homogeneous Markov chain, find the steady state probability of the chain given by the TPM $P = \begin{bmatrix}0.1 & 0.6 & 0.3\\ 0.5 & 0.1 & 0.4\\ 0.1 & 0.2 & 0.7\end{bmatrix}$}{05}{L2}{3}
\q{4b.}{Explain linear and decay model with an example.}{05}{L1}{1}
\q{5.}{A gambler has \rupee 2. He bets \rupee 1 at a time and wins \rupee 1 with probability $\frac{1}{2}$. He stops playing if he ends up with no money or wins \rupee 4.\newline
a. What is the TPM of the related Markov chain?\newline
b. What is the probability that he has lost his money at the end of 5 plays?\newline
c. What is the probability that the game lasts more than 7 plays?}{10}{L3}{4}
\end{longtable}
\end{xpaper}

% ================= BLOCK C: older papers (appended at the end of the file)
% ---------- HS251TA Principles of Management and Economics, Improvement CIE, 27th January 2025
\begin{iempaperB}
\paperinfo{ayear={2024-2025 (ODD Sem)},date={27\textsuperscript{th} January 2025},code={HS251TA},sem={V},exam={Improvement CIE},maxmarks={10 + 50},duration={120 Min},
 title={PRINCIPLES OF MANAGEMENT AND ECONOMICS}}
\iemheadB
\begin{cietable}{M}{BT}{CO}
\qpart{Part -- A}
\q{1}{Which theory compares an individual's input-output ratio to others?}{01}{1}{3}
\q{2}{In Vroom's Expectancy Theory, what does the term \textit{Valence} refer to?}{01}{1}{3}
\q{3}{Mention any one difference between transactional leadership from transformational leadership?}{01}{1}{3}
\q{4}{In Blake and Mouton's Managerial Grid, which leadership style balances concern for people and production?}{01}{1}{2}
\q{5}{According to McGregor's Theory X, what type of behavior is assumed of employees?}{01}{1}{3}
\q{6}{In the IS-LM model, what does the IS curve represent?}{01}{1}{5}
\q{7}{Which Keynesian model incorporates both the goods and labor markets to determine output and employment?}{01}{1}{5}
\q{8}{What is the key assumption of the Keynesian Cross model?}{01}{1}{5}
\q{9}{What role does fiscal policy play in the Keynesian Cross model?}{01}{1}{4}
\q{10}{What does the aggregate supply (AS) curve in the AS-AD model represent?}{01}{1}{5}
\qpart{Part -- B}
\q{1}{A mid-sized company is facing low productivity and high turnover due to employees feeling undervalued and lacking growth opportunities. Using Herzberg's Two-Factor Theory, analyze the hygiene factors and motivators that could improve employee satisfaction.}{10}{4}{3}
\q{2}{A project manager notices that employees' performance drops when they do not understand how their efforts contribute to organizational goals. Using Vroom's Expectancy Theory, explain how the manager can enhance motivation.}{10}{3}{3}
\q{3}{A team leader in a tech company has a diverse team with varying skill levels. Some need detailed guidance, while others prefer autonomy. Using Hersey and Blanchard's Situational Leadership Model, suggest leadership styles for different team members based on their maturity levels.}{10}{3}{3}
\q{4}{Describe the Keynesian Cross Model. How does it explain the equilibrium level of national income in an economy? Discuss the impact of changes in government spending and investment on national income in this model.}{10}{2}{4}
\q{5}{Explain the IS-LM model. Discuss the equilibrium in the goods market (IS curve) and money market (LM curve). How do changes in fiscal and monetary policy affect output and interest rates according to this model?}{10}{2}{5}
\end{cietable}
\btfoot{|l|l|l|*{11}{C{0.8cm}|}}{\hline
\multirow{2}{*}{\shortstack{Marks\\Distribution}} & \multicolumn{2}{c|}{Particulars} & CO1 & CO2 & CO3 & CO4 & CO5 & L1 & L2 & L3 & L4 & L5 & L6\\\cline{2-14}
 & Max & Quiz & -- & 01 & 04 & 01 & 04 & 10 & -- & -- & -- & -- & --\\\cline{3-14}
 & Marks & Test & -- & -- & 30 & 10 & 10 & -- & 20 & 20 & 10 & -- & --\\\hline}
\end{iempaperB}

% ---------- 21IE6F8 Elements of Financial Management, Improvement CIE, 28th August 2024
\begin{iempaperB}
\paperinfo{ayear={2023-2024 (EVEN Sem)},dept={INDUSTRIAL ENGINEERING \& MANAGEMENT},date={28\textsuperscript{th} August 2024},code={21IE6F8},sem={VI},exam={Improvement CIE},maxmarks={10 + 50},duration={20 + 90 min},
 title={Elements of Financial Management (Institutional Elective)}}
\iemheadB
\begin{cietable}{M}{BT}{CO}
\qpart{Part -- A}
\q{1}{Define the term Par Value.}{01}{--}{--}
\q{2}{What is Pre-emptive Right?}{01}{--}{--}
\q{3}{Mention any two phases of the IPO process.}{01}{--}{--}
\q{4}{What is Private Placement?}{01}{--}{--}
\q{5}{In which year the National Stock Exchange was established?}{01}{--}{--}
\q{6}{What is a Merger?}{01}{--}{--}
\q{7}{What is the role of World Bank in the International Financial System?}{01}{--}{--}
\q{8}{What is letter of Credit?}{01}{--}{--}
\q{9}{Define Hedging.}{01}{--}{--}
\q{10}{What type of forward contract is a futures contract?}{01}{--}{--}
\qpart{Part -- B}
\q{1}{Enumerate the principle tasks of Securities Exchange Board of India (SEBI)? Describe briefly the key initiatives taken by SEBI.}{10}{--}{--}
\q{2a}{List the factors that have an important bearing on working capital needs?}{05}{--}{--}
\q{b}{Distinguish between operating cycle and cash cycle?}{05}{--}{--}
\q{3a}{Distinguish between Initial Public Offer (IPO) and Follow on Public Offer (FPO).}{05}{--}{--}
\q{b}{``Accruals are a `free' source of finance''. Analyze.}{05}{--}{--}
\q{4}{Discuss the plausible reasons for Mergers. In what way do you think mergers can enable organizations to beat competitions?}{10}{--}{--}
\q{5}{Illustrate the salient features of foreign exchange market in India. How does it differ from the features of international foreign exchange market?}{10}{--}{--}
\end{cietable}
\end{iempaperB}

\end{document}
